\section{Proof of Proposition~\ref{prop1}}\label{theory:C}

First, we will show that any predecessor of $G$ defines a unique 
sequentially maximal event of $G$. Indeed, for any predecessor $H$ of $G$,
let $e_H$ be the event added from $H$ to reach $G$.
It is immediate to see that $e_H \neq e_H'$ if $H \neq H'$.
Now, let $e = e_H$ for some such predecessor $H$.
Inspecting the definition of $\rightarrow$ in $\mathcal{T}(P)$, it is easy to see that there is no $e'$ in $G.E$ such that $(e,e') \in G.\po{} \cup G.\rf{} \cup G.\mo{}$. Finally, the only event $e'$ for which $(e,e') \in G.\rf{}^{-1}$ could possibly happen is $e' = H.\momax{\loc(e)}$. However,
by definition there cannot be any $e''$ such that $(e',e'') \in G.\mo{}$ and so there cannot be any $e'$ such that $(e,e') \in G.\fr{}$.
Hence $e$ is a sequentially maximal event.

Next, we will show that any sequentially maximal event gives rise to a unique predecessor. Indeed, for any sequentially maximal event $e$ of $G$, let $H_e = G \setminus \{e\}$. Since $e$ is sequentially maximal, removing 
$e$ preserves all relation properties, i.e. $H_e.\rf{} \subseteq G.\rf{}$ and $H_e.\mo{}\subseteq G.\mo{}$. Thus, $H_e$ is a consistent execution graph. Note that $H_e \neq H_{e'}$ if $e \neq e'$. Furthermore $G$ can be obtained from $H_e$ by simply adding $e$ back in. Hence, the only thing we need to show is that $H_e$ is reachable.

To this end, since $H_e$ is sequentially consistent, we can pick a sequentially maximal event $e_1$ of $H_e$ and remove it from $H_e$ to get $H_{e_1}$, which will also be sequentially consistent, by the same argument as before. From there, we can remove another sequentially maximal event and so on and so forth, till we finally reach the initial execution graph $S_0$. This allows us to trace a path from $H_e$ all the way to $S_0$,
which proves that $H_e$ is reachable.